%%%%

%%%   Самарский государственный технический университет

%%%   Кафедра прикладной математики и информатики

%%%

%%%   Преамбула для статьи в журнал "Вестник Самарского государственного

%%%   технического университета. Серия: «Физико-математические науки»"

%%%   Ориентирована под MikTEx 2.4 for Win32

%%%

%%%   Хотя сам журнал собирается под GNU/Linux на дистрибутиве TeXLive



\documentclass[11pt,a4paper,twoside]{article}



\usepackage[cp1251]{inputenc}

\usepackage[T2A]{fontenc}

\usepackage[english,russian]{babel}

\usepackage{samgtubib}

\usepackage[dvips]{graphicx}

\usepackage[multiple]{footmisc}



\usepackage{samgtu}

\renewcommand{\VestnikYear}{2009} % Год издания Вестника

\renewcommand{\VestnikNumber}{2\,(19)} % Номер Вестника

\renewcommand{\VestnikMonth}{Ноябрь} % Месяц выхода Вестника

\renewcommand{\VestnikData}{01.11.09} % Дата подписания Вестника в печать

\renewcommand{\VestnikUPL}{26,64} % Усл. печ. л.

\renewcommand{\VestnikUIL}{25,73} % Уч.-изд. л.

\renewcommand{\VestnikRegNumber}{479/09} % Регистрационный номер

\renewcommand{\VestnikOrder}{994} % Номер заказа







%

%\usepackage{algorithm2e}

%\usepackage{pst-barcode}

%\usepackage{auto-pst-pdf}

%%

%%\newcommand{\totop}[1]{\raisebox{6pt}[1pt][2pt]{#1}} 

%%\newcommand{\tottop}[1]{\raisebox{12pt}[1pt][2pt]{#1}} 

%%\newcommand{\totttop}[1]{\raisebox{18pt}[1pt][2pt]{#1}} 

%%\newcommand{\tottttop}[1]{\raisebox{24pt}[1pt][2pt]{#1}} 

%%\newcommand{\totttttop}[1]{\raisebox{30pt}[1pt][2pt]{#1}} 

%%\newcommand{\tobot}[1]{\raisebox{-6pt}[1pt][2pt]{#1}} 

%%\newcommand{\tobbot}[1]{\raisebox{-12pt}[1pt][2pt]{#1}} 

%%\newcommand{\tobbbot}[1]{\raisebox{-18pt}[1pt][2pt]{#1}}

%

%\graphicspath{{./00PIC/}}

%

%%\samgtupubl % для генерации pdf, отдаваемого в типографию СамГТУ

%\pdfcrop % обрезка pdf для размещения на math-net.ru

%\draft % На корректуре

%\filllastpage % Последняя страница не заполнена не полностью

%%\briefcomm % Статья является кратким сообщением



\vsgtu{1473}



\FirstPage{158}

\LastPage{166}

%%

%%

%\begin{document}



\renewcommand{\baselinestretch}{0.9}



\hfill \qrcode[hyperlink,height=.8in]{http://dx.doi.org/10.14498/vsgtu1437}

\vspace{-.8in}



%\hfill \begin{pspicture}(1in,1in)

%\psbarcode{http://dx.doi.org/10.14498/vsgtu1473}{}{qrcode}

%\end{pspicture}

%\vspace{-1in}







\udc{519.7} \msc{90C39}



\rutitle{К динамическому программированию по~значениям в~полугруппе}

{К динамическому программированию по~значениям в~полугруппе}



\ruauthor{В.~Г.~Овчинников} {О\,в\,ч\,и\,н\,н\,и\,к\,о\,в~В.~Г.}



\ruaffil{\rusamgtu.}



%\email{1}{ (V.~G.~Ovchinnikov,\CAut)}



\ruauthordetails{1}{\fullauthor{\rupid{39274}{Валерий Гаврилович Овчинников}} \regalia{\mem{ovg.samara@mail.ru}} \position{старший преподаватель}\dept{каф. разработки нефтяных и газовых месторождений}}



\enauthordetails{1}{\fullauthor{\enpid{39274}{Valery G. Ovchinnikov}} \regalia{\mem{ovg.samara@mail.ru}} \position{Senior Lecturer}\dept{Dept. of Oil and Gas Fields Development}}



\rukeywords{линейно упорядоченная абелева полугруппа, дискретное

оптимальное управление, оптимальный процесс, доставляющие значения

Беллмана элементы ограничивающих множеств, динамическое

программирование, лексикографические произведения, алгоритмы}





\ruabstract{Для не рассматривавшейся ранее со значениями целевой

функции в линейно упорядоченной абелевой полугруппе $P$ задачи

дискретного оптимального управления даются характеризация

разрешимости и на ее основе алгоритм, ищущий оптимальный процесс,

используя доставляющие значения Беллмана элементы ограничивающих

множеств. Отмечаются модификации данного алгоритма, когда

\begin{itemize}

\item[ $1)$]

$P$ "--- непустое естественно упорядоченное подмножество чисел \linebreak с~операцией получения максимума из двух чисел; 

\item[ $2)$] $P$ "---

естественно упорядоченное множество неотрицательных чисел со

сложением (умножением); 

\item[ $3)$] $P$ "--- лексикографическое

произведение $m$ (не менее двух) линейно упорядоченных абелевых

полугрупп; 

\item[ $4)$] $P$ "--- лексикографическое произведение $m$ (не

менее двух) множеств вещественных чисел с естественным порядком и

сложением, и данный алгоритм получает $m$ "--- оптимальный процесс

проще, чем предыдущий алгоритм автора.

\end{itemize}

}



\entitle{On dynamic programming on the values in~the~semigroup}



\enauthor{V.~G.~Ovchinnikov}{O\,v\,c\,h\,i\,n\,n\,i\,k\,o\,v~V.~G.}



\enaffil{\ensamgtu.}





\enabstract{For not considered previously discrete optimal control problem with target function values in a linearly ordered Abelian semigroup given characterization of the solvability and on its basis the algorithm seeks optimal process with the help of delivering Bellman values elements of limiting sets. We mark the modifications to this algorithm, when

\begin{itemize}

\item[$1)$] $P$ is nonempty subset of numbers with the natural ordering and the operation producing the maximum of two numbers;

\item[$2)$] $P$ is  set of nonnegative numbers with the natural ordering and the addition (or multiplication); 

\item[$3)$] $P$ is lexicographical product of $m$ (not less than two) linearly ordered Abelian semigroups; 

\item[$4)$] $P$ is lexicographic product of $m$ (not less than two) sets of real numbers with the natural ordering and the addition, and this algorithm gets $m$-optimal process easier than the previous author's algorithm.

\end{itemize}

}



\enkeywords{linearly ordered Abelian semigroup, discrete optimal control, optimal process, delivering Bellman values elements of limiting sets, dynamic programming, lexicographical products, algorithms}





\date{04/II/2016} % Дата поступлени  статьи в редакцию.

\revisiondate{22/II/2016} % В окончательном виде

\accepteddate{26/II/2016}



\makerutitle







{\bf 1.} Статья продолжает исследования применения динамического

программирования~\cite{ovchinnikov:bib1,ovchinnikov:bib2}. В ней

для упрощения алгоритма~\cite{ovchinnikov:bib2} используется не

рассматривавшаяся ранее в общей формулировке и в указанных ниже

частных случаях следующая задача дискретного оптимального

управления.



\smallskip



\hypertarget{ovch:zadmm}{}



\begin{task}[$\mathfrak{A}$] Пусть известно следующее\/{\rm :}





\begin{itemize}



\item[--] пространства $A,$ $B;$



\item[--] множество шагов $T(1<| T| <\infty )$,  наделенное

строгим порядком $\prec,$ иными словами$,$  отношением $\prec$ со

свойствами 

$$(i\prec j) \Rightarrow ({i\ne j})(\forall i,j\in

T),\quad  (i\prec g)\wedge (g\prec j)\Rightarrow (i\prec j)(\forall

i,g,j\in T),$$ 

имеющее подмножество начальных шагов 

$$T_0 =\{i\in

T\vert \emptyset =\{j\in T\vert j\prec i\}\}(\ne T),$$ обозначаемые

$\{\pi(i)\}$ одноэлементные подмножества 

$$\big \{j\in T\vert (j\prec

i)\wedge (\emptyset =\{g\in T\vert j\prec g\prec i\})\big\}(\forall i\in

T\backslash T_0 )$$ и удовлетворяющее равенству $\emptyset =M\cap

T_0$ подмножество финальных шагов 

$$M=\{i\in T\vert \tau

(i)=\emptyset\},$$ где $\tau (i)=\{j\in T\backslash T_0 \vert \pi

(j)=i\};$



\item[--] функции $X,$ $U,$ $S$ с конечными множественными значениями 

$$X(i)\subset A(\forall i\in T),\quad U(i,\alpha )\subset B(\forall i\in

T\backslash T_0 )(\forall \alpha \in A),$$ 

 $$S(i,\alpha ,\beta )\subset A(\vert S(i,\alpha ,\beta )\vert \le 1)(\forall i\in T\backslash T_0)

 (\forall \alpha \in A) (\forall \beta \in B);$$



\item[--] обозначение $D$ множества называемых процессами пар $(x,u)$ функций $x,$ $u$ со значениями $x_i \in A(\forall i\in T),$ $u_i \in B(\forall i\in T)$ при ограничениях

\begin{gather*}

x_i =u_i (\forall i\in T_0 ),

\quad

x_i \in X(i)(\forall i\in T),

\quad

u_i \in U(i,x_{\pi (i)} )

\quad

(\forall i\in T\backslash T_0 ),\\

x_i \in S(i,x_{\pi (i)} ,u_i )

\quad

(\forall i\in T\backslash T_0 );

\end{gather*}



\item[--] множество $P$ с операцией $\bot$  и порядком $\trianglelefteq$ как расширением строгого порядка $\triangleleft $ по правилам 

$$(p_1 \trianglelefteq p_2

)\Leftrightarrow (p_1 \triangleleft p_2 )\vee (p_1 =p_2 )(\forall p_1 ,p_2

\in P),$$ являющееся линейно упорядоченной абелевой полугруппой $\langle P ,\bot ,\trianglelefteq \rangle $ и$,$ таким образом$,$ имеющее свойства

\begin{itemize}

\item[] $(p_1 \bot p_2 )\bot p_3 =p_1 \bot (p_2 \bot p_3 ) (\forall p_1 ,p_2 ,p_3 \in P)$ $($ассоциативность $\bot),$



\item[]  $p_1 \bot p_2 =p_2 \bot p_1

(\forall p_1 ,p_2 \in P)$ $($коммутативность $\bot),$



\item[]  

$(p_1 \trianglelefteq p_2 )\vee (p_2 \trianglelefteq p_1)

(\forall p_1 ,p_2 \in P)$ $($линейность $\trianglelefteq),$



\item[]  

$(p_1 \trianglelefteq p_2 )\Rightarrow (p_1 \bot p_3 )\trianglelefteq (p_2 \bot

p_3)(\forall p_1 ,p_2 ,p_3 \in P)$ $($стабильность $\trianglelefteq$ относительно

$\bot);$



\end{itemize}



\item[--] обозначение обобщённой операции $\mathop{\text{\huge$\perp$}}\limits_{i\in I} $ с результатом

$\mathop{\text{\huge$\perp$}}\limits_{i\in I} q(i)$ как элементом $\Lambda (\vert I\vert ,\omega,q)\in P,$ определяемым по

функции $$q : I\to P(\emptyset \ne I\subset T)$$ из условий

$$k>1\Rightarrow \Lambda (k,\omega,q)=\Lambda (k-1,\omega,q) \bot

q(\omega (k))(\forall k\in N_{\vert I\vert } ),$$ 

$$\Lambda

(1,\omega,q)=q(\omega (1))$$ независимо от выбора биекции

$$\omega : N_{\vert I\vert} \to I\quad (N_m =\{1,\ldots,m\}(\forall

m\in \mathbb{N}))$$ $($ср. с замечаниями

N.B.~{\rm \cite[c.~62]{ovchinnikov:bib3}$);$}



\item[--] затраты шагов $$f(i,\gamma ,\alpha ,\beta )(\in P)(\forall i\in T\backslash T_0

)(\forall \gamma ,\alpha \in {\rm A})(\forall \beta \in B);$$



\item[--] обозначение  ${\min \limits_{\gamma \in

\Gamma}}^{(\trianglelefteq)} Q(\gamma )$ наименьшего по

порядку $\trianglelefteq$ значения функции $$Q: \Gamma

\to P.$$





\end{itemize}







Требуется найти называемую оптимальным процессом пару

${(x^\ast,u^\ast)\in D,}$ удовлетворяющую равенству

$$

\mathop{\text{\huge$\perp$}}\limits_{i\in T\backslash T_0} f\left(i,x_i^\ast,x_{\pi

(i)}^\ast,u_i^\ast\right)={\min \limits_{(x,u)\in

D}}^{(\trianglelefteq)} \mathop{\text{\huge$\perp$}}\limits_{i\in

T\backslash T_0} f(i,x_i,x_{\pi (i)},u_i).

$$



\end{task}





\smallskip



\hypertarget{ovch:zam1}{}



\begin{remark}[1]

Если $$P\subset W\in \{\mathbb{R}, \mathbb{Q}, \mathbb{Z}, \mathbb{N}\}(1<\vert P\vert ),\quad   p\bot

q=\max (p,q)(\forall p,q\in P),$$ $\trianglelefteq$

означает $\le$, то оптимальный процесс $(x^\ast,u^\ast)\in D$ находится из условия

\[

\max_{i\in T\backslash T_0} f(i,x_i^\ast,x_{\pi (i)}^\ast, u_i^\ast )=\min \limits_{(x,u)\in D}

\max_{i\in T\backslash T_0} f(i,x_i,x_{\pi (i)},u_i

).

\]

\end{remark}



\smallskip



\hypertarget{ovch:zam2}{}



\begin{remark}[2]

Если 

$$P=\{p\in W\vert 0\le p\}(W\in \{\mathbb{R}, \mathbb{Q}, \mathbb{Z}, \mathbb{N}\}),$$  $$p\bot

q=p+q \; \text{$($либо $p\bot q=p\cdot q)$ $(\forall p,q\in P)$,}$$ 

$\trianglelefteq$ означает $\le$, то оптимальный процесс $(x^\ast,u^\ast)\in D$ находится из условия

\begin{gather*}

\sum _{i\in T\backslash T_0 } f(i,x_i^\ast,x_{\pi

(i)}^\ast,u_i^\ast )=\min \limits_{(x,u)\in D}

\sum _{i\in T\backslash T_0 } f(i,x_i,x_{\pi (i)},u_i)\\

\biggl(\text{либо} \quad \prod _{i\in T\backslash T_0 } f(i,x_i^\ast,x_{\pi (i)}^\ast ,u_i^\ast)=\min

\limits_{(x,u)\in D} \prod _{i\in T\backslash T_0}

f(i,x_i,x_{\pi (i)},u_i) \biggr).

\end{gather*}

\end{remark}



\smallskip



\hypertarget{ovch:zam3}{}



\begin{remark}[3] В задаче \hyperlink{ovch:zadmm}{$\mathfrak{A}$} линейно упорядоченная абелева полугруппа $\langle

P,\bot,\trianglelefteq \rangle $ может получаться как

лексикографическое произведение линейно упорядоченных абелевых

полугрупп $\langle P_k,\bot _k,\trianglelefteq_k \rangle

(\forall k\in N_m )(m\in \mathbb{N}\backslash \{1\})$ или, иными

словами, определяться из следующих условий:



\begin{itemize}

\item[--]

$P$ "--- множество функций $$q:  N_m \longrightarrow \bigcup

\limits_{k\in N_m} P_k$$ со значениями $q_k \in P_k$ ($ \forall k

\in N_m),$



\item[--]

$(p\trianglelefteq q)\Leftrightarrow(p=q)\vee (\emptyset

\ne K=\{k\in N_m \vert p_k \ne q_k \})\wedge (e={\min K}) \wedge

(p_e\trianglelefteq_e q_e)$ $(\forall p,q\in P)$,



\item[--]

$(p\bot q)_k =p_k \bot _k q_k $ $(\forall k\in N_m )$ $(\forall

p,q\in P)$.





\end{itemize}

\end{remark}



\smallskip



\hypertarget{ovch:zam4}{}



\begin{remark}[4] Если полугруппа $\langle P,\bot, \trianglelefteq

\rangle $ определяется по замечанию \hyperlink{ovch:zam3}{3}, 

$P_k =\mathbb{R}(\forall k\in N_m )$, $\trianglelefteq_k $ означает $\le (\forall k\in N_m )$, $\bot

_k $ означает $+$ $(\forall k\in N_m )$ и, кроме того, 

$$S(i,\alpha,\beta )=\{s(i,\alpha,\beta)\}$$

и 

$$\big(f(i,s(i,\alpha,\beta),\alpha,\beta)\big)_k=f^k(i,\alpha,\beta) (\forall i\in T\backslash

T_0)  (\forall \alpha \in A)(\forall \beta \in B)(\forall k\in N_m ),$$

то оптимальный процесс $(x^\ast,u^\ast)\in D$

является $m$-оптимальным процессом~\cite{ovchinnikov:bib2}.

\end{remark}



\smallskip



{\bf 2.} В общем случае наряду с обозначениями~\cite{ovchinnikov:bib2} 

$$T_k =\{i\in T\vert

h(i)=k\}(\forall k\in N_{h(T)}),$$ 

где $$h(T)=\max \limits_{i\in T} h(i),\quad h(i)=\vert \{j\in T\vert j\prec i\}\vert,$$

будет использоваться следующее развитие терминов~\cite{ovchinnikov:bib2}: 



\begin{itemize}

\item[--]

пара

$(Y,V)$ называется {\sl подходящей}, если её составляют функции

$Y$, $V$ со значениями 

$$Y(i)\subset A(\forall i\in T), \;  V(i,\alpha )\subset B(\forall \alpha \in Y(\pi (i))) (\forall i\in

T\backslash T_0),$$ когда $Y(i)\ne \emptyset (\forall i\in

T)$; 

\item[--]

для подходящей пары $(Y,V)$ множество

$F(Y,V)$ образуют пары $(x,u)$ функций $x$, $u$

со значениями $$x_i \in A(\forall i\in T),\quad u_i \in B(\forall

i\in T)$$ 

при ограничениях 

$$x_i =u_i (\forall i\in T_0 ), \quad x_i \in

S(i,x_{\pi(i)},u_i )(\forall i\in T\backslash T_0

),$$ $$x_i \in Y(i)(\forall i\in T),\quad u_i

\in V(i,x_{\pi(i\mbox)})(\forall i\in T\backslash T_0

);$$ 

\item[--]

подходящая пара $(Y,V)$ называется

$S$-{\sl согласованной}, если выполнены условия 

$$\emptyset \ne V(i,\alpha )=\{\beta \in

V(i,\alpha )\vert \emptyset \ne S(i,\alpha ,\beta )\cap Y(i)\} (\forall

\alpha \in Y(\pi (i))) (\forall i\in T\backslash T_0).$$



\end{itemize}



\smallskip



{\bf 3.} Следующее утверждение (теорема~\hyperlink{ovch:th1}{1}), развивающее теорему~1~\cite{ovchinnikov:bib2},

даёт характеризацию (необходимые и достаточные условия) разрешимости задачи~\hyperlink{ovch:zadmm}{$\mathfrak{A}$}.



\smallskip



\hypertarget{ovch:th1}{}



\begin{theorem}[1]Оптимальный процесс задачи \hyperlink{ovch:zadmm}{$\mathfrak{A}$} существует$,$ если и только если из равенств

$$X^\partial (i)=X(i)(\forall i\in M),$$ $$

X^\partial (i)=\{\alpha \in X(i)\vert \emptyset \ne \{\beta \in

U(j,\alpha)\vert \emptyset \ne S(j,\alpha ,\beta)\cap

X^\partial (j)\}(\forall j\in \tau (i))\}$$ $$

(\forall i\in T_{k-1} \backslash M) (\forall k\in N_{h(T)})$$ $($по уменьшению $k)$ получаются непустые множества достижимости 

$$X^\partial (i)(\forall i\in T).$$ 

При этих условиях определяемая значениями $$X^\partial (i)\subset A (\forall i\in T)$$ функция

$X^\partial $ и определяемая значениями 

$$U^\partial (i,\alpha )=\{\beta \in U(i,\alpha )\vert \emptyset

\ne S(i,\alpha ,\beta )\cap X^\partial (i)\}\subset B(\forall \alpha \in

X^\partial (\pi (i)))(\forall i\in T\backslash T_0 )$$ функция $U^\partial$ составляют $S$"=согласованную пару $(X^\partial ,U^\partial ),$ удовлетворяющую равенству $D=F(X^\partial ,U^\partial).$

\end{theorem}



\smallskip



{\bf 4.} Далее предполагается, что оптимальный процесс задачи \hyperlink{ovch:zadmm}{$\mathfrak{A}$}

существует, пара $(X^\partial,U^\partial )$ составлена, как указано в

теореме~\hyperlink{ovch:th1}{1}, и, следовательно, 

$$

D=F(X^\partial,U^\partial),\quad 

\emptyset \ne X^\partial (i)(\forall i\in T\mbox),$$ 

$$S(i,\alpha ,\beta)\ne

\emptyset \ne U^\partial (i,\alpha )=\{\beta \in U^\partial (i,\alpha )\vert

\emptyset \ne S(i,\alpha ,\beta )\cap X^\partial (i)\}

$$ $$(\forall \beta \in

U^\partial (i,\alpha ))(\forall \alpha \in X^\partial (\pi

(i)))(\forall i\in T\backslash T_0).$$



Также предполагается, что так называемые значения Беллмана ${\bf

B}(i,\alpha )$ определены из условий

\begin{gather*}

(i\in M)\Rightarrow {\bf B}(i,\alpha )=\min_{\beta \in

U^\partial (i,\alpha)}^{{\hspace{9mm}}(\trianglelefteq)}

f(i,\psi(S(i,\alpha,\beta)),\alpha ,\beta),\\

(i\notin M)\Rightarrow {\bf B}(i,\alpha )= \min \limits_{\beta \in

U^\partial (i,\alpha )}^{{\hspace{9mm}}(\trianglelefteq)}

\Bigl(f(i,\psi (S(i,\alpha ,\beta)),\alpha ,\beta)\bot

\mathop{\text{\huge$\perp$}}\limits_{j\in \tau (i)}

{\bf B}(j,\psi (S(i,\alpha,\beta))) \Bigr)\\

(\forall \alpha \in X^\partial (\pi (i))

(\forall i\in T_k )

\left(\forall k\in N_{h(T)} \right) \quad \text{(индукцией по уменьшению} \,\, k),

\end{gather*}

где $\psi$ "--- функция выбора  при условиях 

$$\emptyset \ne C\Rightarrow \psi

(C)\in C(\forall C\subset A\cup B\cup T).$$



Следующее утверждение (теорема~\hyperlink{ovch:th2}{2}) в определённых пределах характеризует

оптимальность процесса задачи \hyperlink{ovch:zadmm}{$\mathfrak{A}$}.



\smallskip



\hypertarget{ovch:th2}{}



\begin{theorem}[2]Процесс $(x^\ast,u^\ast) \in D$ задачи~\hyperlink{ovch:zadmm}{$\mathfrak{A}$} оптимален$,$

если $($и при стабильности $\triangleleft$ относительно $\bot, $ только если$)$ одновременно выполнены условия

\begin{gather*}

{\bf B}\left(i,x_{\pi (i)}^\ast \right)=f\left(i,x_i^\ast,x_{\pi

(i)}^\ast ,u_i^\ast \right) \quad

(\forall i\in M), \\

{\bf B}(i,x_{\pi (i)}^\ast )=f\left(i,x_i^\ast,x_{\pi (i)}^\ast

,u_i^\ast \right)\bot \mathop{\text{\huge$\perp$}}\limits_{j\in \tau (i)}

{\bf B}\left(j,x_i^\ast \right) \quad

(\forall i\in (T\backslash T_0 )\backslash M),\\

\mathop{\text{\huge$\perp$}}\limits_{j\in \tau (i)} {\bf B}(j,x_i^\ast )= \min

\limits_{\alpha \in X^\partial (i)}^{{\hspace{9mm}}(\trianglelefteq)}

\mathop{\text{\huge$\perp$}}\limits_{j\in \tau (i)} {\bf B}(j,\alpha) (\forall i\in

T_0 ).

\end{gather*}

\end{theorem}



\smallskip



\begin{remark}[5] В случае замечания~\hyperlink{ovch:zam1}{1} (когда нет стабильности $\triangleleft $

относительно $\bot$) легко строится задача~\hyperlink{ovch:zadmm}{$\mathfrak{A}$} с оптимальным процессом, не

удовлетворяющим всем условиям теоремы~\hyperlink{ovch:th2}{2}.

\end{remark}



\smallskip



Из теоремы~\hyperlink{ovch:th2}{2} получается



\smallskip



\begin{newthm}{Следствие}

Значения функций, составляющих оптимальный процесс \linebreak $(x^\ast,u^\ast)\in D$ задачи~\hyperlink{ovch:zadmm}{$\mathfrak{A}$}, находит следующий алгоритм динамического программирования $($алгоритм \hyperlink{ovch:alg}{$\mathfrak{B}$}$).$

\end{newthm}



\smallskip



\hypertarget{ovch:alg}{}



\begin{newthm}{Алгоритм $\mathfrak{B}$} Этот алгоритм образуют описываемые при помощи псевдоязыка \verb"Pascal" {\rm \cite{ovchinnikov:bib4}} следующие последовательно выполняемые действия \linebreak \rm $(${\tt Цикл~1}, {\tt Цикл 2}, {\tt Цикл 3}$)$:

\end{newthm}



\smallskip



\hypertarget{ovch:loop1}{}



{\tt Цикл 1}, находящий элементы ${\mathbf x}(i,\alpha )(\in S(i,\alpha

,{\mathbf u}(i,\alpha)))$, ${\mathbf u}(i,\alpha )(\in U^\partial

(i,\alpha ))$ из условий

\begin{gather*}

i\in M\Rightarrow {\bf B}(i,\alpha )=f(i,{\mathbf x}(i,\alpha

),\alpha

,{\mathbf u}(i,\alpha)),\\

i\notin M\Rightarrow {\bf B}(i,\alpha )=f(i,{\mathbf x}(i,\alpha

),\alpha ,{\mathbf u}(i,\alpha ))\bot \mathop{\text{\huge$\perp$}}\limits_{j\in

\tau (i)}

{\bf B}(j,{\mathbf x}(i,\alpha ))\\

(\forall \alpha \in X^\partial (\pi (i))

(\forall i\in T\backslash T_0 ):

\end{gather*}



\hspace{3mm}{\bf for} $k:=h(T)$ {\bf downto} 1 {\bf do}



\hspace{3mm}{\bf for}  для каждого $i\in

T_k$ {\bf do}



\hspace{3mm}{\bf for} для каждого $\alpha \in X^\partial (\pi

(i))$ {\bf do}



\hspace{3mm}{\bf begin} $\beta^\ast:=\psi (U^\partial

(i,\alpha))$; $r^\ast :=f(i,\psi (S(i,\alpha,\beta^\ast)),

       \alpha,\beta^\ast)$;



\hspace{16mm}{\bf if} $\tau (i)\ne \emptyset$ {\bf then}



\hspace{16mm}{\bf for} для

каждого $j\in \tau (i)$ {\bf do} $r^\ast :=r^\ast \bot {\bf B}(j,\psi

(S(i,\alpha,\beta^\ast)))$;



\hspace{16mm}{\bf if} $U^\partial (i,\alpha )\backslash \{\beta^

\ast \}\ne \emptyset$ {\bf then}



\hspace{16mm}{\bf  for} для каждого $\beta \in

U^\partial (i,\alpha )\backslash \{\beta^\ast\}$ {\bf do}



\hspace{16mm} {\bf begin} $r:=f(i,\psi

(S(i,\alpha,\beta)),\alpha,\beta)$;



\hspace{27mm}{\bf if} $\tau (i)\ne \emptyset$ {\bf then}



\hspace{27mm}{\bf for} для каждого $j\in \tau (i)$

{\bf do} $r:=r\bot {\bf B}(j,\psi(S(i,\alpha,\beta)))$;



\hspace{27mm}$b:=${\bf not}$(r^\ast \trianglelefteq r)$;

{\bf if} $b$ {\bf then begin} $r^\ast :=r$;

$\beta^\ast :=\beta$ {\bf end};



\hspace{16mm}{\bf end};



\hspace{16mm}${\bf B}(i,\alpha ):=r^\ast$;

$\mathbf{x}(i,\alpha):=\psi(S(i,\alpha,\beta^\ast))$;

$\mathbf{u}(i,\alpha ):=\beta^\ast$



\hspace{3mm}{\bf end}.



\smallskip



\hypertarget{ovch:loop2}{}



{\tt Цикл 2}, $(\forall i\in T_0 )$ получающий $x_i^\ast$, $u_i^\ast$,

используя $$\mathbf{B}(j,\alpha) (\forall \alpha \in X^\partial (\pi

(j)))(\forall j\in T_1):$$



\hspace{3mm}{\bf for} для каждого $i\in T_0$ {\bf do}



\hspace{3mm}{\bf begin} $\Delta :=\tau (i)\backslash \{\psi (\tau (i))\}$;

$\alpha ^\ast :=\psi (X^\partial (i))$;

$r^\ast :={\bf B}(\psi (\tau (i)), \alpha ^\ast )$;



\hspace{13mm} {\bf if} $\Delta \ne \emptyset$ {\bf then}



\hspace{13mm} {\bf for} для каждого $j\in \Delta$ {\bf do}

$r^\ast:=r^\ast \bot {\bf B}(j,\alpha^\ast)$;



\hspace{13mm} {\bf if} $X^\partial (i)\backslash \{\alpha ^\ast

\}\ne \emptyset$ {\bf then}



\hspace{13mm} {\bf for} для каждого $\alpha \in X^\partial

(i)\backslash \{\alpha ^\ast \}$ {\bf do}



\hspace{13mm} {\bf begin} $r:={\bf B}(\psi (\tau (i)),\alpha)$;



\hspace{23mm} {\bf if} $\Delta \ne \emptyset$ {\bf then}



\hspace{23mm} {\bf for} для каждого $j\in \Delta$ {\bf do}

$r:=r\bot {\bf B}(j,\alpha)$;



\hspace{23mm} $b:={\bf not}(r^\ast \trianglelefteq r)$;

{\bf if} $b$ {\bf then}

{\bf begin} $r^\ast :=r$; $\alpha ^\ast :=\alpha$

{\bf end}



\hspace{13mm} {\bf end};



\hspace{13mm} $u_i^\ast :=x_i^\ast :=\alpha ^\ast$



\hspace{3mm} {\bf end}.



\smallskip



{\tt Цикл 3}, $(\forall i\in T\backslash T_0 )$ получающий $x_i^\ast$,

$u_i^\ast $  с использованием $x_i^\ast (\forall i\in T_0 )$,

найденных циклом \hyperlink{ovch:loop2}{2}, и использованием ($\forall \alpha \in

X^\partial (\pi (i))$ $(\forall i\in T\backslash T_0)$ элементов

${\bf x}(i,\alpha)$, $\mathbf{u}(i,\alpha)$, найденных циклом \hyperlink{ovch:loop1}{1}:



\hspace{6mm} {\bf for} $k:=1$ {\bf to} $h(T)$ {\bf do for}  для каждого $i\in

T_k$ {\bf do}



\hspace{6mm} {\bf begin}



\hspace{6mm} $x_i^\ast :={\bf x}(i,x_{\pi(i)}^\ast )$; $u_i^\ast

:={\bf u}(i,x_{\pi (i)}^\ast )$

{\bf end}.



\smallskip



\begin{remark}[6] В случае замечания~\hyperlink{ovch:zam1}{1}  алгоритм \hyperlink{ovch:alg}{$\mathfrak{B}$} модифицируется по правилам

\begin{gather*}

(b:={\bf not}(p\trianglelefteq q)) \Leftrightarrow

(b:={\bf not}(p\le q))

(\forall p,q\in P),\\

(p:=p\bot q) \Leftrightarrow

(p:=\max (p,q))

(\forall p,q\in P).

\end{gather*}

\end{remark}



\smallskip



\begin{remark}[7] В случае замечания~\hyperlink{ovch:zam2}{2}, алгоритм \hyperlink{ovch:alg}{$\mathfrak{B}$} модифицируется по правилам

\begin{gather*}

(b:={\bf not}(p\trianglelefteq q)) \Leftrightarrow

(b:={\bf not}(p\le q)) (\forall p,q\in P),\\

 (p:=p\bot q) \Leftrightarrow

(p:=p+q(p:=p\cdot q))

(\forall p\mbox{,}q\in P).

\end{gather*}

\end{remark}



\smallskip



\hypertarget{ovch:zam8}{}



\begin{remark}[8] В случае замечания~\hyperlink{ovch:zam3}{3} алгоритм \hyperlink{ovch:alg}{$\mathfrak{B}$} модифицируется по правилам

\begin{center}

$(b:={\bf not}(p\trianglelefteq q)) \Leftrightarrow $



$(b:=$ {\bf false}; $k:=1$; {\bf while} $(p_k =q_k)\wedge (k\leq m)$  {\bf do} $k:=k+1;$



{\bf if} $k\le m$ {\bf then} $b:={\bf not}( p_k

\trianglelefteq_k q_k )) (\forall p,q\in P)$,



$(p:=q)\Leftrightarrow$ ({\bf for} $k:=1$ {\bf to} $m$ {\bf do}

$p_k :=q_k$)($\forall p,q\in P)$,



$(p:=p\bot q)\Leftrightarrow$ ({\bf for} $k:=1$ {\bf to} $m$ {\bf

do} $p_k :=p_k \bot_k q_k$) $(\forall p,q\in P)$.

\end{center}

\end{remark}



\smallskip



\begin{remark}[9] В случае замечания~\hyperlink{ovch:zam4}{4} модифицированный по правилам замечания~\hyperlink{ovch:zam8}{8}

алгоритм \hyperlink{ovch:alg}{$\mathfrak{B}$} получает $m$-оптимальный процесс проще алгоритма \cite{ovchinnikov:bib2} за счет

неиспользования, нехранения и неполучения $s$-согласованных пар

$(X^k,U^k)(\forall k\in N_m )$, даваемых теоремой 2~\cite{ovchinnikov:bib2}.



\end{remark}



\smallskip



{

\thanks{{\bf ORCID} 



{\bf Валерий Гаврилович Овчинников:} \url{http://orcid.org/0000-0003-2234-5010}



}}





\newpage



\RusRef





\begin{thebibliography}{9}



\RBibitem{ovchinnikov:bib1}

\by Овчинников~В.~Г.

\paper Алгоритмы динамического программирования оптимальных и близких к~ним процессов

\inbook Труды пятой Всероссийской научной конференции с~международным участием \rm (29--31 мая 2008~г.). Часть~4

\bookinfo Информационные технологии в~математическом моделировании

\serial Матем. моделирование и краев. задачи

\yr 2008

\pages 107--112

\publ СамГТУ

\publaddr Самара

\mathnet{http://mi.mathnet.ru/mmkz1174}



\RBibitem{ovchinnikov:bib2}

\by Овчинников~В.~Г.

\paper К алгоритмам динамического программирования оптимальных процессов

\jour Вестн. Сам. гос. техн. ун-та. Сер. Физ.-мат. науки

\yr 2012

\issue 3(28)

\pages 215--218

\mathnet{http://mi.mathnet.ru/vsgtu1102}

\crossref{10.14498/vsgtu1102}



\RBibitem{ovchinnikov:bib3}

\by Фор~Р., Кофман~А., Дени--Папен~М.

\book Современная математика

\publaddr М.

\publ Мир

\yr 1966

\totalpages 271



\RBibitem{ovchinnikov:bib4}

\by Ахо~А.~В., Хопкрофт~Д.~Э., Ульман~Д.~Д.

\book Структуры данных и алгоритмы

\publaddr М.

\publ Вильямс

\yr 2009

\totalpages 400



\end{thebibliography}



\RuDate





\bigskip



\makeentitle



{

\thanks{{\bf ORCID} 



{\bf Valery G. Ovchinnikov:} \url{http://orcid.org/0000-0003-2234-5010}



}}







\EngRef





\begin{thebibliography}{9}



\Bibitem{ovchinnikov:bib1:eng}

\lang In Russian

\by Ovchinnikov~V.~G.

\paper Algorithms of dynamic programming for optimal and similar processes

\inbook Proceedings of the Fifth All-Russian Scientific Conference with international participation \rm (29--31 May 2008). Part~4

\serial Matem. Mod. Kraev. Zadachi

\yr 2008

\pages 107--112

\publ Samara State Technical Univ.

\publaddr Samara





\Bibitem{ovchinnikov:bib2:eng}

\lang In Russian

\by Ovchinnikov~V.~G.

\paper On the algorithms of dynamic programming for optimal processes

\jour Vestn. Samar. Gos. Tekhn. Univ. Ser. Fiz.-Mat. Nauki

\yr 2012

\vol 3(28)

\pages 215--218

\crossref{10.14498/vsgtu1102}



\Bibitem{ovchinnikov:bib3:eng}

\by Faure~R., Kaufmann~A., Denis-Papin~M.

\book Manuale di Matematica

\publ ISEDI

\publaddr Milano 

\yr 1975

\lang In Italian



\Bibitem{ovchinnikov:bib4:eng}

\by Aho~A.~V., Hopcroft~J.~E., Ullman~J.~D.,

\book  Data Structures and Algorithms

\publ Addison-Wesley

\yr 1983

\serial Addison-Wesley Series in Computer Science and Information Processing

\publaddr Reading, Massachusetts, etc.

\totalpages xi+427 

\zmath{https://zbmath.org/?q=an:0487.68005}





\end{thebibliography}



\EngDate



\newpage

\Russian





%\end{document}

